\documentclass[11pt]{article}

\usepackage[review]{acl}
\usepackage{times}
\usepackage{latexsym}
\usepackage[T1]{fontenc}
\usepackage[utf8]{inputenc}
\usepackage{booktabs}
\usepackage{microtype}

\title{From Concern Detection to Promotion Control:\\
A Protocol for Fail-Closed Governance of Autonomous Research Artifacts}
\author{Anonymous}
\date{}

\begin{document}
\maketitle

\begin{abstract}
Autonomous research systems increasingly produce complete-looking manuscripts
and run bundles, yet review components can identify serious scientific defects
without preventing a paper-ready decision. We formulate this gap as
\emph{promotion governance}: deciding whether an artifact bundle should be
promoted, reviewed, downgraded, or blocked, while linking each blocking concern
to inspectable evidence and a responsible repair stage. We introduce a
manifest-driven evaluation protocol with clean controls, counterfactual fault
variants, base-bundle-disjoint splits, blind manuscript-only requests, and
metrics for false promotion, concern--acceptance conflict, clean-case retention,
repair localization, traceability, and cost. A synthetic development suite
with four base bundles and 40 cases validates the evaluator and exposes three
missing integrity checks in an artifact audit. After those checks were added,
the audit passed all development cases; an advisory-only ablation still
accepted every defective case despite detecting its concerns. These results
are instrument validation, not empirical evidence for general effectiveness:
the suite is synthetic, unreviewed, and explicitly ineligible for paper-level
claims. We specify an oracle-certified controlled experiment with held-out
fault families and repeated real-model evaluation needed to test the central
hypothesis without treating human adjudication as metric gold.
\end{abstract}

\section{Introduction}

Research-agent systems span literature review, experimentation, and report
writing \cite{schmidgall-etal-2025-agent}. Capability benchmarks now evaluate
replication workflows extending from paper understanding to experiment
execution \cite{starace2025paperbench}. Completion, however, is not equivalent
to scientific readiness. BadScientist shows that plausible but unsound papers
can receive favorable automatic reviews
\cite{jiang-etal-2026-badscientist}, while counterfactual evaluation finds that
flaws in research logic have no significant effect on automatic reviews
\cite{dycke-gurevych-2026-automatic}.

We study the policy boundary between detecting a concern and allowing an
artifact bundle to advance. A useful research governor must fail closed: a
blocking concern cannot coexist with a paper-ready decision unless the concern
is resolved or explicitly downgraded. It must also avoid the trivial solution
of blocking everything. This motivates three questions: whether artifact-bound
gates reduce false promotion, whether binding concerns to transitions prevents
concern--acceptance conflict, and what clean-case, repair-localization, latency,
and cost trade-offs follow.

The contribution at the present evidence level is a protocol and tested
instrumentation, not a confirmed performance claim. We provide (i) a portable
artifact-bundle task with explicit promotion and repair labels, (ii) leakage
and independent-oracle gates, (iii) baselines that isolate detection from action binding,
and (iv) a development failure analysis that routes audit misses back to the
responsible workflow stages.

\section{Related Work}

\paragraph{Research-agent capability.}
Agent Laboratory organizes literature review, experimentation, and reporting
within a multi-agent research workflow \cite{schmidgall-etal-2025-agent}.
PaperBench evaluates replication of published machine-learning research from
paper understanding through executed artifacts \cite{starace2025paperbench}.
These systems motivate artifact-level evaluation, but their primary object is
research capability rather than the consistency of a promotion decision with
detected blockers.

\paragraph{Review and claim verification.}
BadScientist operationalizes concern--acceptance conflict for unsound generated
papers \cite{jiang-etal-2026-badscientist}. Dycke and Gurevych use targeted
edits to create flawed research logic and compare automatic reviews
\cite{dycke-gurevych-2026-automatic}. CLAIM-BENCH evaluates claim extraction,
evidence extraction, and claim--evidence link validation in AI papers
\cite{javaji-etal-2025-ai}. These dimensions motivate a separate evaluation
question: whether detected concerns constrain promotion decisions, bind artifact evidence, and identify the responsible repair stage.

\paragraph{Artifact assurance and recovery.}
ARA reconstructs scientific workflows for reproducibility assessment
\cite{riehl2026ara}. REPRO-Bench evaluates agents' reproducibility assessments
using paper PDFs and reproduction packages \cite{hu2025reprobench}. SAGE
routes experiment failures to structured intervention levels
\cite{ma2026reflection}. The closest distinction is therefore not
artifact inspection or failure attribution alone, but whether detected
scientific blockers deterministically constrain paper-readiness promotion
without over-blocking clean bundles.

\section{Promotion Governance Protocol}

\subsection{Task and labels}

Each case is a portable research-run artifact tree paired with a private case
manifest. The evaluated system returns one decision from \texttt{promote},
\texttt{needs\_review}, \texttt{downgrade}, or \texttt{block}; zero or more
concerns with severity and evidence paths; and the workflow stages responsible
for repair. Gold labels remain outside the visible artifact tree.

Every suite declares an evidence class, evaluation regime, claim ceiling,
adjudication status, mutation-isolation status, and paper-claim eligibility.
The controlled regime derives gold from a frozen fault registry and requires
an independently implemented oracle to replay each mutation and verify the
resulting artifact hash. Its claim ceiling is limited to registered fault
families and it cannot claim external validation. A separate naturalistic
regime retains independent human adjudication and external source provenance.
All variants derived from one base bundle remain in the same split, and source
content hashes detect renamed duplicates across splits.

\subsection{Fault families and controls}

The protocol requires clean positive, null, and negative controls together
with variants covering missing comparison evidence, repeated-run provenance,
hidden failed execution, declared-versus-executed budget mismatch,
result--figure conflict, claim--evidence conflict, citation support mismatch,
stale persisted state, and unsupported claim strength. Each mutation records
before and after hashes. For controlled evaluation, the implementation assigns
entire fault families to development or test, derives case gold from the
registry, and verifies every clean and mutated artifact by independent replay.
Certification binds registry, gold, split, oracle-report, and development-suite
hashes; any mismatch quarantines the suite. Evaluated systems receive only the
artifact view, and blind manuscript requests omit case, family, path, oracle,
and gold metadata. Naturalistic evaluation can additionally use separate
mutation auditors and promotion-label reviewers.

\subsection{Comparison conditions}

We define four deterministic conditions and one provider condition.
\emph{Always promote} trusts the incoming ready state. A \emph{presence
checklist} verifies required files without semantic consistency checks. An
\emph{advisory artifact audit} reports the same concerns, evidence, and repair
owners as the full audit but does not bind them to the decision. The
\emph{fail-closed artifact audit} applies the binding policy. Finally, a blind
\emph{manuscript-only reviewer} receives only paper text under opaque request
identifiers; case labels, mutations, paths, and gold data remain private.

\subsection{Metrics and inference}

Primary metrics are false paper-ready promotion and concern--acceptance
conflict. We also report four-way decision macro-F1, clean-case promotion,
blocker precision/recall/F1, exact repair-owner match, trace coverage, latency,
and provider cost. Every system trial must cover every case. Pairwise effects
are computed per case, with 5,000-replicate bootstrap intervals clustered by
base bundle and exact paired sign tests over base-bundle effects. Synthetic
development suites remain exploratory regardless of numeric performance.

\section{Implementation and Development Validation}

The protocol is implemented as manifest-driven builders, runners, scorers, and
blind request import/export commands. Paths are confined to suite roots;
symbolic links, duplicate cases, hash mismatches, split leakage, and incomplete
trial coverage are validation failures. Failure analysis emits review
diagnostics and node-strengthening recommendations with explicit recheck
conditions.

\begin{table}[!t]
\centering
\small
\begin{tabular}{lrrr}
\toprule
Condition & False & Conflict & Clean \\
\midrule
Always promote & 1.00 & -- & 1.00 \\
Presence checklist & 1.00 & -- & 1.00 \\
Advisory audit & 1.00 & 1.00 & 1.00 \\
Fail-closed audit & .00 & .00 & 1.00 \\
\bottomrule
\end{tabular}
\caption{Synthetic development results over four base bundles and 40 cases.
Columns report false-promotion, concern--acceptance conflict, and clean-case
promotion rates.
The labels were generated for instrument debugging and were not independently
adjudicated; the table is not confirmatory evidence.}
\label{tab:development}
\end{table}

The development corpus contains four synthetic base bundles, one clean control
and nine counterfactual variants per bundle, for 40 cases. It is marked as
synthetic development evidence with unreviewed adjudication and no paper-claim
eligibility. Table~\ref{tab:development} therefore tests instrumentation only.

Before strengthening, the fail-closed audit reached .70 decision accuracy and
.333 false-promotion rate, missing repeated-run provenance, executed-budget,
and stale-state faults. The failure analyzer routed the first two to experiment
execution and the third to review. Added integrity checks were then evaluated
on the same frozen suite. In the final development run, the fail-closed audit
differed from the advisory ablation by .90 decision accuracy and -1.00 false
promotion. The cluster bootstrap intervals are degenerate because each
synthetic fault repeats identically across only four base bundles; the paired
sign-test value is .125. This is direct evidence that the development sample is
too small for inferential claims, despite final decision accuracy and blocker
F1 of 1.00 on the generated cases.

\section{Confirmatory Study and Readiness Gate}

The controlled confirmatory study requires at least 720 held-out cases. Entire
registered fault families are assigned to either development or test before
outcomes are observed; base identities and source hashes are also disjoint.
Every test base contributes one clean control and one case from each held-out
test family. Certification independently replays all development and test
mutations from their clean controls, verifies registry-derived gold, and binds
the complete case set and split manifests by hash. A certified suite may become
paper-claim eligible for registered fault families without human label
adjudication, but it remains explicitly unvalidated on naturally occurring
defects.

The manuscript-only condition requires three preserved real-model trials per
case. Remote API trials retain hashed provider receipts; local trials retain an
exact model-artifact digest and hash-bound runtime receipts. Deterministic
baselines and ablations use versioned, hash-bound system-run manifests. At
least one post-repair rerun per original fault case must measure recovery and
clean-control regression. Raw requests, responses, decisions, failures,
manifests, hashes, latency, and cost must be retained. The confirmatory gate
checks held-out family coverage for every system and trial, recomputes metrics
from predictions, and cannot raise the claim ceiling beyond the registered
families.

The central hypothesis is that fail-closed governance lowers false promotion
by at least 20 absolute percentage points relative to the presence checklist
while retaining at least 90 percent clean-case promotion. The study is stopped
or downgraded if the full policy blocks more than 10 percent of clean bundles,
if the checklist is within five points on both primary safety and clean-case
metrics, if mutations leak labels, if oracle replay or split isolation fails,
or if required real-model and recovery evidence is incomplete.

\section{Limitations and Ethical Considerations}

Counterfactual faults may be easier than naturally occurring research defects,
and deterministic checks can overfit a known artifact schema. A promotion gate
does not establish scientific truth; it enforces consistency between available
evidence, stated concerns, and allowed transitions. Provider-based reviewers
introduce model, prompt, and temporal dependence. Public artifact bundles must
exclude credentials, private paths, personal data, and proprietary research
content. False promotion can amplify unsound claims, while false blocking can
suppress valid null or negative results; both errors must remain visible.

\section{Conclusion}

Promotion governance treats paper readiness as an evidence-bound transition
rather than a by-product of workflow completion. The implementation and
development suite show that this distinction is measurable and that detected
failures can drive targeted audit strengthening. They do not yet show that the
policy generalizes. An oracle-certified controlled suite, repeated real-model
review, complete recovery evidence, and a frozen confirmatory run are required
before advancing beyond a research memo. Naturalistic generalization remains a
separate claim that would require independently curated external evidence.

\bibliography{references}

\end{document}
